\scanpage{47}
\mspage{22}
% archive sheet number: 46

COROLLAIRE.- \uline{Soit} $\handwritten{u}(z_1,..,z_p)$ \uline{une fonction holomorphe dans une partie ouverte} $A$ \uline{de} $K^p$, \uline{à valeurs dans} $K^p$. \uline{Pour tout point} $\handwritten{a}=(a_1,..,a_p)$ \uline{de} $A$ \uline{tel que la différentielle} $\handwritten{\unsure{d.u(a)|du(a)}}$ \uline{soit inversible il existe un voisinage ouvert} $U$ \uline{de} $\handwritten{a}$ \uline{tel que la restriction de} $\handwritten{u}$ \uline{à} $U$ \uline{soit un homéomorphisme de} $U$ \uline{sur un voisinage ouvert} $V$ \uline{de} $\handwritten{u(a)}$, \uline{et que l'homéomorphisme réciproque soit holomorphe dans} $V$.

\margin{\handwritten{\figure{\begin{tikzpicture}[x=0.10cm,y=0.10cm,baseline=0pt]
\fill (0,0) -- (-0.15,1.8) -- (-1.15,1.5) -- (-0.8,0.4) -- cycle;
\draw[line width=0.55pt] (-0.4,1.35) -- (-1.6,1.1) -- (-2.9,0.9);
\draw[line width=0.45pt] (-0.55,0.6) .. controls (-1.3,1.05) and (-1.8,0.2) .. (-2.15,0.55);
\end{tikzpicture}}}}
Lorsque $K=\handwritten{\unsure{\mathbb{R}|R}}$ et que $A$ est un intervalle ouvert dans $\handwritten{R}$, le fait que $u'(x)\handwritten{\ne}0$ dans $A$ entraîne que $u$ est un homéomorphisme de $A$ \uline{tout entier} sur $u(A)$ puisque $u$ est monotone dans $A$ (\uline{Fonct.var.réelle}, chap.I,\S\,2 et \uline{Top.gén.} chap.IV,\S\,2). Mais lorsque $K=\handwritten{\mathbb{C}}$, $u'(x)$ peut être $\handwritten{\ne}0$ en tout point d'un domaine $A\handwritten{\subset C}$ sans que $u$ soit un homéomorphisme de $A$ sur $u(A)$. Par exemple, si $u(x)=x^n$, on a $u'(x)\handwritten{\ne}0$ pour $x\handwritten{\ne}0$, mais si $A$ est le disque $\handwritten{|x-1|<\frac{1}{2}}$, $u$ n'est pas un homéomorphisme de $A$ sur $u(A)$ dès que $n$ est assez grand, car il existe des points $x\handwritten{\ne}1$ dans $A$ tels que $x^n=1$.

Il faudrait faire ici également le théorème de \uline{dépendance} des fonctions holomorphes à jacobien identiquement nul.

\uline{3}. \uline{Intégration des équations différentielles}.

THEOREME 2.- \uline{Soit} $K$ \uline{un corps valué complet de caractéristique} $0$. \uline{Soient} \margin{\handwritten{\unsure{-}}}$f_i(x;y_1,..,y_p;z_1,..,z_q)$ $(1\handwritten{\leq}i\handwritten{\leq}p)$ $p$ \uline{fonctions holomorphes scalaires définies dans une partie ouverte} $A$ \uline{de} $K^{p+q+1}$. \uline{Pour tout point} $(a;b_1,..,b_p;c_1,..,c_q)$ \uline{de} $A$, \uline{il existe un voisinage} $U$ \uline{de} $a$ \uline{dans} $K$, \uline{un voisinage} $V$ \uline{de} $(b_1,..,b_p)$ \uline{et un voisinage} $W$ \uline{de} $(c_1,..,c_q)$ \uline{tels que} $U\handwritten{\times}V\handwritten{\times}W\handwritten{\subset}A$, \uline{et un système et un seul de} $p$ \uline{fonctions scalaires} $u_i(x;z_1,..,z_q)$, \uline{holomorphes dans} $U\handwritten{\times}W$ \uline{telles que l'image de} $U\handwritten{\times}W$ \uline{par l'application} $(x;z_1,..,z_q)\handwritten{\to}(u_i(x;z_1,..,z_q))$ \uline{soit contenue dans} $V$ \uline{et qui satisfont dans} $U\handwritten{\times}W$ \uline{aux conditions}
